\documentclass{article}
\usepackage{amsmath,amsfonts,amssymb}
\usepackage{tikz}
\usepackage{tikz-cd}
\usepackage{tikz}
\usepackage{float}
\numberwithin{figure}{section}
\tikzset{
vertex/.style={
draw,
circle,
fill=gray!20,
minimum size=20pt,
inner sep=1pt
}
}
\title{Thesis}
\author{Derrickson Ymbon}
\begin{document}
\maketitle
\section{Introduction: Preliminaries}

\subsection{Introduction}
The representation theory of finite-dimensional algebras has been a very active area in the world of mathematics as of late. More importantly, there is a trichotomy in the theory, where the representation type of algebras can be of three types namely, finite, tame or wild. Algebras of finite type are most easily understood. An algebra is said to be of finite type if the number of isoclasses of indecomposable modules is finite. Algebras of tame representation type, on the other hand, have an infinite number of isoclasses of indecomposable modules, but our understanding of them is more-or-less ``controlled", in the sense that all but a finite number of these indecomposable modules can be parametrized by 1-parameter families, and wild representation type is the least understood among the three.\\ 

There is an important class of algebras called the \textit{special biserial} algebas that was shown to be tame\cite{Wald1985TameBA, waschbusch1983representation}. This class contains some other classes of algebras which are crucial in representation theory, but here we will mainly be interested in two, namely the gentle algebras and Brauer graph algebras\cite{schroll2017brauergraphalgebras}.\\
Gentle algebras and Brauer graph algebras are closely related, that is, we can obtain Brauer graph algebras using the trivial extension of the gentle algebras, and we can go back to the original algebra via the concept of admissible cuts.\cite{schroll2015trivialextensionsgentlealgebras} \\

We will also look at the triangulations of a marked oriented surface $S$ whose marked points lie on the boundary of $S$, and the algebras that arise from such triangulations, and thus investigate the connections between those algebras and the algebras that we already know of\cite{assem2009gentlealgebrasarisingsurface}.\\ 

If we have a quiver $Q$ with vertex set $Q_0$, arrow set $Q_1$ and if we denote the start of an arrow $\alpha$ as $s(\alpha)$ and its endpoint as $t(\alpha)$, we can obtain the corresponding \textit{path algebra} $KQ$, whose underlying vector space has the basis being the set of all paths in $Q$. If $I$ is an admissible ideal of $KQ$, then $KQ/I$ is a finite dimensional algebra called a bound quiver algebra, with $(Q,I)$ called a bound quiver\cite{schiffler2014quiver,assem2006elements}. In this study, we will be interested in the path algebras where certain ``constraints/restrictions" are imposed on the quiver. An example of such a restriction is that there is at most one arrow leaving and entering every vertex in the quiver. However, there will be 4 main restrictions that we will apply. They are as follows: 

\begin{enumerate}
    \item[G1)] For every vertex $v\in Q_0$, there are at most two arrows entering $v$ and at most two arrows leaving $v$
    \item[G2)] For every arrow $\beta$ in $Q$, there is at most one arrow $\alpha$ such that $\alpha\beta \notin I$ and at most one arrow $\gamma$ such that $\beta\gamma\notin I$.
    \item[G3)] For every arrow $\beta$ in $Q$, there is at most one arrow $\alpha$ such that $\alpha\beta \in I$ and at most one arrow $\gamma$ such that $\beta\gamma \notin I$.
    \item[G4)] The ideal $I$ is generated by paths of length 2 
\end{enumerate}

Conditions G1 and G2 give us a nice symmetric and compact(not in the topological sense) view of the quiver, and as one would expect, the path algebra obtained from this type of quiver naturally has interesting properties too. This type of algebra is known as the special biserial algebra. Further, on imposing conditions G3 and G4, we get a gentle algebra, which has even more intricacies to its structure. \\

Here, there is the notion of a trivial extension of an algebra, whereby we get a new algebra  by increasing the dimension of the underlying vector space of the original algebra. Formally, the trivial extension of the $K-$algebra $A$ is defined to be $T(A)=A\ltimes D(A)$, where $D(A)=\mathrm{Hom}_K(A,K)$. This is the algebra whose underlying $K-$vector space is given by $A \oplus D(A)$, and where the multipilication defined by:
\[
(a,\phi)(b,\psi)=(ab,a\psi+\phi b)
\]

Interestingly, this algebra happens to be a Brauer graph algebra, and as such there must be a Brauer graph and thus a quiver associated to this algebra.

There will also be an exposition on special types of graphs called \textit{ribbon} graphs or \textit{fat} graphs. These graphs are special because they have a geometric interpretation or realisation. That is, given a ribbon graph $\Gamma$, we have the geometric realisation of the graph to be the topological space
\[
|\Gamma|=(E \times [0,1])/ \sim
\]

where $\sim$ is the equivalence relation generated by $(e,t) \sim (\iota(e), 1-t)$ for all $t\in [0,1]$ and $(e,0) \sim (f,0)$ if $s(e)=s(f)$ and $(e,1)\sim (f,1)$ if $s(\iota(e)=s(\iota(f)$, for all $e,f \in E$. Here $E$ is the set of half-edges in $\Gamma$ and $s$ and $\iota$ are maps that sends a half edge to the vertex it is incident on and a half-edge incident on a vertex to the half-edge incident on the other vertex of the same edge respectively.\\

This ribbon graph can also be considered as a ``generalisation" of a Brauer graph, where the cyclic ordering in the Brauer graph can be taken to be the permutation of the edges in the ribbon graph, and further the vertices of the ribbon graph are weighted to give the multiplicity of the Brauer graph. \\

Further, ribbon graphs are significant because they can be embedded into an open oriented surface with boundary where the cyclic orderings around the vertices of the graph is induced by the orientation of the surface. This statement unleashes a lot of capabilities of ribbon graphs, and there are multiple theorems which depend on the above, such as the fact that the embeddings of two ribbon graphs into two different surfaces gives us a uniqueness of the surfaces upto homeomorphism, provided that the graphs are isomorphic.  \\

Given a gentle algebra $A$, we can also obtain its ribbon graph by a certain technique. This puts even more importance on the ribbon graphs and adds an extra layer in the interplay between the different structures and the corresponding fields of geometry, combinatorics and algebra.\\

Finally, given a triangulation $T$ of a marked surface $(S,M)$, we can obtain a quiver $Q(T)$, the path algebra of which will be a gentle algebra. More explicitly, first we consider an oriented surface $S$ with no punctures and whose marked points $M$ all lie in the boundary of $S$,i.e., $M\subset \partial S$. Then if we have a triangulation $T$ of the surface, more specifically an ideal triangulation, then we associate to $T$ a quiver $Q(T)$. From this quiver, we can get its corresponding path algebra $A_T$, which is a gentle algebra(!!).\\

So, ultimately, we have fascinating connections between the different structures presented in this survey, that is, we have back-and-forth connections between gentle algebras and Brauer graph algebras, between gentle algebras and triangulations, between gentle algebras and ribbon graphs, and finally between  ribbon graphs and oriented surfaces. This provides a lot of insight into the interaction of the different fields of 
geometry, algebra, combinatorics and even a slight touch of topology.
\subsection{Preliminaries: K-Algebras, Modules}

The representation theory of finite-dimensional algebras can be seen as the study of modules over these algebras. So, in this section, we will begin by defining objects called \textit{modules} and \textit{K-algebras}.\\


First we begin with a survery of special types of rings called $K$-\textit{algebras}.\\

\textbf{Definition 1.2.1: } \textit{Let $K$ be an algebraically closed field. A $K$-algebra is a ring $A$ with identity element 1 together with a vector space structure that is compatible with the ring multiplication inside $A$. That is, there is an operation $\cdot : K \times A \to A, (\lambda,a) \mapsto \lambda a$, such that} 
$$
\lambda (ab)=(\lambda a)b=(a\lambda)b=a(\lambda b)=(ab)\lambda
$$

\textbf{Examples 1.2.1: } \textit{
    \begin{itemize}
        \item The polynomial ring $K[x]$ in the indeterminate $x$ is a $K$-algebra. In fact it is an infinite dimensional $K$-algebra.
        \item The ring $M_n(K)$ of all square matrices of order $n$ with entries from $K$ is an $n^2$-dimensional $K$-algebra
    \end{itemize}
}
Now we define the radical of an algebra, which is an important object.\\

\textbf{Definition 1.2.2: } \textit{Let $K$ be an algebraically closed field. The radical of a $K$-algebra $A$, also sometimes called the Jacobson radical, denoted by rad $A$ is defined as the intersection of all maximal right ideals of $A$. That is, if $M$ is the set of all maximal right ideals in $A$, then}
\[
\text{rad } R= \bigcap_{I\in M } I
\]

\textbf{Examples 1.2.2: } \textit{
    \begin{enumerate}
        \item If we consider a field $K$ as an algebra over itself, then rad $K$ is simply the zero ideal. That is rad $K$=0
        \item Let A be the $K$-algebra $K[x]/(x^s)$. Then the Jacobson radical of A is given by the ideal $(x+(x^s))$
    \end{enumerate}
}
\textbf{Defintion 1.2.3: } \textit{
    Let $R$ be a ring, and let $A$ be an $R$-algebra. Then we have the following definitions:
    \begin{enumerate}
        \item An element $e \in A$ is called an \textbf{idempotent} element if $e^2=e$.
        \item An element $e$ is said to be a \textbf{central} idempotent if it is an element of the center of $A$,i.e. $ea=ae \forall a \in A$
        \item Two idempotents $e_1$ and $e_2$ are called \textbf{orthogonal} if $e_1e_2=e_2e_1=0$.
        \item An idempotent is said to be {primitive} if $e$ cannot be decomposed into two non-zero orthogonal idempotents. That is, if $e=e_1+e_2$ for some orthogonal elements $e_1,e_2$, then either $e_1=0$ or $e_2=0$  
    \end{enumerate}
}

\textbf{Definition}

\textbf{Definition }: \textit{Let $R$ be a ring with 1. A (left) R-module is a set M equipped with a binary operation $+: M \times M \to M$ such that $(M,+)$ is an abelian group, and an action of the ring $R$ on $M$(i.e a map $R \times M \to M$ denoted by $r\cdot m,\forall r\in R, m \in M$) such that $\forall r,s \in R, m\in M$, the following properties are satisfied}:

\begin{itemize}
    \item[(M1)] $(r+s)\cdot m=r\cdot m + s\cdot m$ 
    \item[(M2)] $r\cdot (m+n)=r\cdot m+r\cdot n$
    \item[(M3)] $(r\cdot s)\cdot m=r\cdot(s\cdot m)$  
    \item[(M4)] $1 \cdot m=m$ 
\end{itemize}

Studying the ring $R$ and its structure becomes a lot easier if we study its modules. \\

\textbf{Examples : } \textit{
    \begin{enumerate}
        \item The ring $R$ can be seen as an R-module with the action being just the ordinary multiplication inside the ring itself.
        \item If $I$ is a left(right) ideal of the ring, then $I$ is a left(right) module over $R$
        \item If $G$ is an abelian group, then $G$ is a $\mathbb{Z}$-module
    \end{enumerate}
}



\section{Special Biserial Algebras and Gentle Algebras}
\subsection{Quivers and Path Algebras}
\textbf{(Defn): }Let $K$ be an algebraically closed field. A quiver is a quadruple $Q=(Q_0,Q_1,s,t)$, where \begin{itemize}
    \item $Q_0$ is the set of all vertices 
    \item $Q_1$ is the set of all arrows
    \item A map $s:Q_1 \to Q_0$ which map an arrow $\alpha \in Q_1$ to its starting point $s(\alpha)$ 
    \item A map $t:Q_1 \to Q_0$, which maps the same arrow to its ending point $t(\alpha)$. 
\end{itemize}
In other words, if $\alpha$ is an arrow, then $s(\alpha)$ is the ``head" of the arrow and $t(\alpha)$ is the ``tail" of the arrow. So, a quiver is essentially just a directed graph where loops and parallel arrows are allowed. As such, when we take a loop $\alpha$, then the source and target of the loop is the same. For example, take the following quiver
\begin{figure}[H]
\centering
\begin{tikzcd}[sep=normal]
	&&& 3 \\
	1 && 2 \\
	&&& 4
	\arrow["\mu", from=1-4, to=1-4, loop, in=55, out=125, distance=10mm]
	\arrow["\alpha", from=2-1, to=2-3]
	\arrow["\beta", shift left=2, from=2-3, to=1-4]
	\arrow["\eta"', shift right=2, from=2-3, to=1-4]
	\arrow["\gamma"', shift right=2, from=2-3, to=3-4]
	\arrow["\lambda"', shift right=2, from=3-4, to=2-3]
\end{tikzcd}
\caption{Quiver $Q^{(1)}$}
\label{fig:quiverexample}
\end{figure}

 The above diagram is a quiver where $Q_0=\{1,2,3,4\},Q_1=\{\alpha,\beta,\gamma,\lambda,\mu,\eta\}$. where $s(\alpha)=1$ and $t(\alpha)=2$,$s(\beta)=2=s(\eta)$ and $t(\beta)=3=t(\eta)$, $s(\mu)=t(\mu)=3$, $s(\lambda)=4$ and $t(\lambda)=2$, and $s(\gamma)=2$ and $t(\gamma)=4$\\


 A \textit{path} in $Q$ from a vertex $i$ to another vertex $j$ is a sequence of arrows  $c=(i\mid\alpha_1,\alpha_2,...\alpha_m\mid j)$, where the target of the previous arrow is the source of the next arrow. That is, in $c$, we have $t(\alpha_{k-1})=s(\alpha_k)$ for all $0\leq k \leq t-1$ . Here, we see $s(\alpha_1)=i$ and $t(\alpha_{m})=j$ and we can also denote the start and ending of the paths by $s(p)=\alpha_1$ and $t(p)=\alpha_m$ respectively. The length of the path $c$ is the number of arrows in the path, denoted by $l$ and thus $l(c)=m$. The \textit{lazy} path or trivial path is simply the path $e_i=(i\mid \alpha \mid i)$ for some arrow $\alpha$. In Figure~\ref{fig:quiverexample}, $c_1=(1\mid \alpha,\beta\mid 2)$ is a path. We can also denote paths just by the arrows and dropping the vertices. Then $c_2=\alpha\gamma\lambda$ is also a path from 1 to 2\\
 

 For every quiver $Q$, we can also get its corresponding \textit{path algebra} $KQ$, which is an algebra over the field $K$, and whose underlying vector space is generated by all the paths in the quiver. In other words, the basis of $KQ$ is the set of all paths in $Q$. Also important to note that if the quiver has a loop or an oriented cycle, then there would be infinite paths, and thus the path algebra would be infinite dimensional. The path algebra for the quiver $Q^{(1)}$ in Figure~\ref{fig:quiverexample} is infinite dimensional whereas the path algebra for a quiver with 2 vertices and a single arrow between them would be finite dimensional\\
 
 Like any other algebra (or vector space), the multiplication can be determined by the multiplication of two basis elements. So, here if $c$ and $c'$ are any two basis elements of $KQ$, then the multiplication $c \cdot c'$ is defined by 

 \[
c \cdot c' = 
\begin{cases}c  c' , & \text{if } s(\alpha_k)=t(\alpha_{k-1}) \\
0, & \text{otherwise }
\end{cases}
\]

where $c c'$ is the concatenation of two paths in $Q$(Note that concatenation of two paths is only possible if the ending vertex of $c$ is the starting vertex of $c'$, so that means multiplication in $KQ$ is possible only if concatenation in $Q$ is possible. If it is not we simply define it to be 0). \\


Let us denote by $\mathcal{R}$ the two-sided ideal of $KQ$ generated by all the arrows in $Q$. Another two sided ideal $I$ is called admissible if there exists a positive integer $m\geq 2$ such that 
\[
\mathcal{R}^m \subset I \subset \mathcal{R}^2
\]

Here, $\mathcal{R}^m$ is simply the decomposition $$\mathcal{R}^m=\bigoplus_{l \geq m}KQ_l$$ and as a vector space, it has a basis being the set of all paths greater than or equal to $m$. So, if $I$ is an admissible ideal, then it would contain all the paths greater than or equal to $m$, which would mean that the algebra obtained by quotienting out $I$ from $KQ$, that is, the algebra $KQ/I$ is finite dimensional. We also say that the algebra $KQ/I$ is given by a quiver and relations. \\
Also, the condition $I \subset\mathcal{R}^2$ guarantees that when we quotient $I$ out from the path algebra, then we would not cut out any arrows.\\

\subsection{Special Biserial Algebras and Gentle Algebras}

As mentioned earlier, a restriction we can impose on the bound quiver $(Q,I)$ is the number of arrows going in and out of a vertex in $Q$. In particular, let us consider those quivers where there are at most two arrows starting and ending at every vertex in the quiver. We will call this condition G1. One more condition is a condition on the ideal where multiple continuations and pre-continuations are killed in the algebra. Formally stated, we have the two conditions:
\begin{enumerate}
    \item[G1)] The number of arrows starting at $v$ and ending at $v$ is at most 2, for every vertex $v$ in the quiver $Q$.
    \item[G2)] For every arrow $\beta$ in $Q$, there exists at most one arrow $\alpha$ such that $\alpha\beta\notin I$ and there exists at most one arrow $\gamma$ such that $\beta\gamma \notin I$.
\end{enumerate}

By placing the conditions above on the quiver, we obtain a special biserial algebra.\\

\textbf{Definition: } A finite-dimensional $K$-algebra $A$ is called \textit{special biserial} if there exists a quiver $Q$ and an admissible ideal $I$ such that $A$ is Morita equivalent to $KQ/I$ 
The conditions are:
\begin{enumerate}
    
    \item[G3)] The ideal $I$ is generated by paths of length 2.
    \item[G4)] For every arrow $\beta$ in $Q$, there exists at most one arrow $\alpha$ such that $\alpha\beta\in I$ and at most one arrow $\gamma$ such that $\beta\gamma \in I$.       
\end{enumerate}
 

 \textbf{Definition: } A finite dimensional algebra $\Lambda$ is called a \textit{gentle} algebra if it is Morita equivalent to the bound quiver algebra $KQ/I$ that satisfies the conditions above. 

 The conditions G1 and G2 are relatively easy to understand. For G2, we can phrase it in another way in that if for every arrow $\beta$ there exist multiple arrows, say $\alpha_1, \alpha_2,\alpha_3$ before $\beta$ in $Q$, then only one of of those continuations with $\beta$ will not be in the ideal,i.e., say $\alpha_1\beta\notin I$, then $\alpha_2\beta\in I$ and $\alpha_3\beta\in I$, and the same goes for arrows after $\beta$. In other words, the paths $\alpha_2\beta$ and $\alpha_3\beta$ will not survive in the algebra. The reasoning is similar for G2. If we take the same example as for G2, then $\alpha_1\beta \in I$, whereas $\alpha_2\beta, \alpha_3\beta\notin I$
 
\section{Geometric Interpretation of Brauer Graph Algebras}

\textbf{Definition 2.1: } A \textit{Brauer graph} $G$ is a tuple $G=(V,E,m,o)$ where 
\begin{itemize}
    \item $(V,E)$ is an unoriented connected graph with vertex set $V$ and edge set $E$.
    \item $m: V \to \mathbb{Z}^+$ is a function, called the \textit{multiplicity}, which assigns to each vertex a positive integer.
    \item $o$ is called the \textit{orientation} of $G$, and is given as follows: for each vertex $v \in V$, there is a cyclic ordering of the edges incident with $v$. For example, this cyclic ordering may be clockwise or anti-clockwise. If $v$ is incident to only one edge $e$ and $m(v)=1$, then the cyclic ordering is simply given by $e$. If $m(v) > 1$, then the cyclic ordering at $v$ is given by $e<e$.
\end{itemize}
%squiggly line between gentle and brauer graph algebras
\section{Surface Triangulations and Connections to Gentle Algebras}
\section{Conclusion and Scope of Work}


\end{document}